\documentclass[10pt,twocolumn]{ctexart}

\usepackage[a4paper,margin=0.7in]{geometry}
\usepackage{lmodern}
\usepackage{microtype}
\usepackage{hyperref}
\usepackage{booktabs}
\usepackage{array}
\usepackage{enumitem}
\usepackage{xcolor}

\hypersetup{
  colorlinks=true,
  linkcolor=black,
  urlcolor=blue,
  pdftitle={火车票务管理系统 - 统一技术规范},
  pdfauthor={项目团队}
}

\setlist[itemize]{leftmargin=*,itemsep=2pt,topsep=2pt}
\setlist[enumerate]{leftmargin=*,itemsep=2pt,topsep=2pt}
\setlength{\columnsep}{0.22in}
\setlength{\parindent}{0pt}
\setlength{\parskip}{4pt}

\title{\textbf{火车票务管理系统}\\统一技术规范}
\author{}
\date{版本 1.0}

\begin{document}
\maketitle
\thispagestyle{empty}

\section{项目概述}
火车票务管理系统是一个基于 Qt Widgets 和 C++ 开发的桌面应用程序。项目采用分层架构，以分离界面展示逻辑、业务逻辑和数据库访问逻辑。

本文档是团队范围内的技术基线，用于定义统一的软件架构、模块职责、数据库设计、接口约定、集成规则和开发约束。

所有成员都应按照本规范完成各自模块的实现。若需要修改本文档或共享架构，应在实施前与项目负责人沟通并达成一致。

\section{技术栈}
本项目使用 \textbf{C++17} 作为编程语言，采用 \textbf{Qt Widgets} 作为桌面界面框架，使用 \textbf{CMake} 进行构建，数据库统一选用 \textbf{SQLite}。版本控制工具为 \textbf{Git + GitHub}，日常开发以 \textit{develop} 为主分支，稳定发布以 \textit{main} 为目标分支。

\section{项目结构}
仓库应保持当前的高层目录结构：

\begin{verbatim}
TrainTicketSystem/
src/
  login/
  train/
  ticket/
  statistics/
  database/
  ui/
include/
database/
docs/
tests/
\end{verbatim}

当前约定的核心类命名也属于共享结构的一部分：
\textit{LoginManager}、\textit{TrainManager}、\textit{TicketManager}、
\textit{StatisticsManager}、\textit{DatabaseManager}。

\section{软件架构}
项目采用如下分层架构：

\begin{verbatim}
Qt UI
  |
  v
LoginManager
TrainManager
TicketManager
StatisticsManager
  |
  v
DatabaseManager
  |
  v
SQLite
\end{verbatim}

\subsection*{架构规则}

UI 类不得直接执行 SQL，所有数据库访问必须统一通过 \textit{DatabaseManager} 完成。各个 Manager 类通过公共接口向 UI 提供服务，并负责各自模块内部的业务逻辑处理。不同模块之间不应直接操作彼此的内部数据；若需要调整共享结构、重命名核心类或改变分层方式，必须先经过讨论再实施。

\section{模块职责}
\subsection*{项目负责人}
项目负责人主要负责数据库设计与实现、GitHub 仓库管理、系统集成、代码审查以及 Pull Request 审查，同时承担团队协作过程中的接口协调和整体推进工作。

\subsection*{程序员 A}
模块：\textit{login/}

程序员 A 负责登录窗口、身份验证逻辑、管理员相关功能以及权限管理实现。

\subsection*{程序员 B}
模块：\textit{train/}

程序员 B 负责站点 CRUD、车次 CRUD、车次查询以及席位库存管理等与列车信息相关的功能。

\subsection*{程序员 C}
模块：\textit{ticket/}、\textit{statistics/}

程序员 C 负责订票、退票、改签、订单历史以及统计分析等票务与报表相关功能。

\subsection*{测试人员}
测试人员负责功能测试、文档编写、用户手册整理以及 Bug 报告提交。

\section{数据库规范}
\subsection*{数据库引擎}：SQLite

\textbf{版本 1 表}

版本 1 的正式表包括 \textit{User}、\textit{Station}、\textit{Train} 和 \textit{Order}。

\textbf{后续版本可扩展表}

若后续需求扩大，可进一步引入 \textit{Ticket}、\textit{Payment}、\textit{Seat} 和 \textit{OperationLog} 等扩展表。

\textbf{版本 1 设计决策}

在首个集成版本中，不单独设置 \textit{Ticket} 表。当前约定为：\textit{Order} 表中的一条记录代表一条购票记录。这样可以简化数据库结构，降低多人协作时的集成风险。若后续需求强制要求座位级建模，可在版本 2 中扩展设计。

\subsection*{User}
\begin{tabular}{>{\raggedright\arraybackslash}p{0.22\columnwidth} >{\raggedright\arraybackslash}p{0.18\columnwidth} >{\raggedright\arraybackslash}p{0.48\columnwidth}}
\toprule
字段 & 类型 & 说明 \\
\midrule
\textit{userId} & INTEGER & 主键，自增 \\
\textit{username} & TEXT & 唯一，非空 \\
\textit{password} & TEXT & 非空 \\
\textit{role} & INTEGER & 0 游客，1 售票员，2 管理员 \\
\textit{enabled} & INTEGER & 默认值为 1 \\
\bottomrule
\end{tabular}

\subsection*{Station}
\begin{tabular}{>{\raggedright\arraybackslash}p{0.24\columnwidth} >{\raggedright\arraybackslash}p{0.20\columnwidth} >{\raggedright\arraybackslash}p{0.44\columnwidth}}
\toprule
字段 & 类型 & 说明 \\
\midrule
\textit{stationId} & INTEGER & 主键，自增 \\
\textit{stationName} & TEXT & 唯一，非空 \\
\bottomrule
\end{tabular}

\subsection*{Train}
\begin{tabular}{>{\raggedright\arraybackslash}p{0.30\columnwidth} >{\raggedright\arraybackslash}p{0.18\columnwidth} >{\raggedright\arraybackslash}p{0.34\columnwidth}}
\toprule
字段 & 类型 & 说明 \\
\midrule
\textit{trainId} & INTEGER & 主键 \\
\textit{trainNumber} & TEXT & 唯一，非空 \\
\textit{departureStationId} & INTEGER & 外键，指向 Station \\
\textit{arrivalStationId} & INTEGER & 外键，指向 Station \\
\textit{departureTime} & TEXT & 非空 \\
\textit{arrivalTime} & TEXT & 非空 \\
\textit{totalSeats} & INTEGER & 非空 \\
\textit{remainingSeats} & INTEGER & 非空 \\
\textit{enabled} & INTEGER & 默认值为 1 \\
\bottomrule
\end{tabular}

必要约束：
\begin{itemize}
  \item \textit{remainingSeats >= 0}
  \item \textit{totalSeats >= 0}
  \item \textit{remainingSeats <= totalSeats}
\end{itemize}

\subsection*{Order}
\begin{tabular}{>{\raggedright\arraybackslash}p{0.25\columnwidth} >{\raggedright\arraybackslash}p{0.18\columnwidth} >{\raggedright\arraybackslash}p{0.41\columnwidth}}
\toprule
字段 & 类型 & 说明 \\
\midrule
\textit{orderId} & INTEGER & 主键 \\
\textit{userId} & INTEGER & 外键，指向 User \\
\textit{trainId} & INTEGER & 外键，指向 Train \\
\textit{passengerName} & TEXT & 非空 \\
\textit{purchaseTime} & TEXT & 非空 \\
\textit{status} & INTEGER & 0 已预订，1 已退票，2 已改签 \\
\bottomrule
\end{tabular}

\subsection*{关系定义}
\textit{Order.userId} 关联 \textit{User.userId}，\textit{Order.trainId} 关联 \textit{Train.trainId}。同时，\textit{Train.departureStationId} 和 \textit{Train.arrivalStationId} 分别关联 \textit{Station.stationId}，用于表示车次的出发站和到达站。

\section{接口约定}
当前共享的核心 Manager 类如下：

\textit{LoginManager}、\textit{TrainManager}、\textit{TicketManager}、\textit{StatisticsManager} 以及 \textit{DatabaseManager}。

各 Manager 通过公共接口向 UI 提供服务。\textit{DatabaseManager} 是唯一允许执行 SQL 的组件。

\subsection*{预期职责}

\textit{LoginManager} 负责登录验证、角色校验和密码修改；\textit{TrainManager} 负责车次与站点的增删改查、车次搜索以及库存更新；\textit{TicketManager} 负责订票、退票、改签和历史查询；\textit{StatisticsManager} 负责报表与统计汇总；\textit{DatabaseManager} 则负责连接管理、数据库模式初始化以及共享查询执行。

\section{开发规则}
每个功能都应在独立的功能分支中开发，所有合并必须通过 Pull Request 完成，禁止直接向 \textit{main} 分支推送。日常开发的主要目标分支为 \textit{develop}。同时，Commit 信息必须清晰描述本次实现内容，便于审查与追踪。

\textbf{良好示例}：\textit{Implement login validation}、\textit{Add train search}、\textit{Fix refund bug}。

\textbf{避免使用}：\textit{update}、\textit{fix}、\textit{123}。

\subsection*{AI 使用规则}

如果团队成员使用 AI 生成代码，则该 AI 的输出也必须遵循本规范，保持既有结构一致，不得擅自发明新的整体架构，也不得绕开既定模块边界。

\section{集成规则}
默认情况下，每位程序员只修改自己负责的模块。若需要跨模块修改，应提前沟通并说明原因。数据库结构变更必须经过项目负责人批准，公共接口变更也应在实施前同步说明。无论模块如何扩展，UI 代码都必须保持不直接包含 SQL 逻辑。

\section{当前项目状态}
\begin{tabular}{>{\raggedright\arraybackslash}p{0.48\columnwidth} >{\raggedright\arraybackslash}p{0.38\columnwidth}}
\toprule
项目项 & 状态 \\
\midrule
架构设计 & 已确定 \\
目录结构 & 已确定 \\
模块职责 & 已确定 \\
数据库设计 & 版本 1 已冻结 \\
数据库实现 & 进行中 \\
SQL 初始化 & 尚未实现 \\
Qt UI & 进行中 \\
测试 & 尚未开始 \\
\bottomrule
\end{tabular}

\end{document}
